\subsubsection{问题三的结论}\label{sec:q3-5}

把本节的做法归纳成对题目的回答：

\begin{itemize}
\item \textbf{搜索策略。}用半径 $999$ m 圆周上等角分布的 $7$ 个检测站保证发现：每站的 $1000$ m 接收圆合起来盖住整个场地，任何位置的全向源都会在某个站被收到，因此不必在原点先扫全部频道。每到一站把未发现频道扫一遍，顺便对附近的已知源补测；扫过几站后，多余的站剪掉、未去的站往路线上挪，每次调整都重新核验覆盖。
\item \textbf{定位策略。}每个已发现的源维护一个凸多边形定位区域，示向度交楔形，无信号点挖去 $1000$ m 圆盘并与有信号点配对切半平面。继续测向时在包围圆圆心附近布设候选点，按移动时间加测后剩余代价选点，最多 $4$ 次后回圆心测，每次至少把半径减半。
\item \textbf{清除策略。}包围圆半径不超过 $19.5$ m 时，走到离区域所有顶点都不超过 $20$ m、且兼顾下一任务的位置清除，一定成功。半径不超过 $80$ m、$20$ m 圆盘能盖住区域面积四成以上时，先试一次光学清除，每源最多两次，失败只花 $3$ s 且不改变区域。只有收到真实的成功反馈才算清除。
\item \textbf{任务安排。}未去的站和已知源合成一个路线问题，用多起点局部搜索排序，只执行第一个任务后重新规划；$16$ 个频道全部清除，或站扫完且已知源全部清除时结束。
\end{itemize}

在本地复原模拟器的 $2000$ 个场景上，这套策略全部完成清除，平均定位清除时间 $239.3$ s/源，比不做提前光学清除的方案省 $1.63$ s/源，在 $1736$ 个场景上更快。它省时间的来源是用便宜的光学尝试替代近距离的重复测向；不漏源、不裁掉真实源、判定能清除时一定能清除这三条，不依赖任何尝试的成败。
